\chapter*{Bibliografia}
\addcontentsline{toc}{chapter}{Bibliografia}
\markboth{Bibliografia}{Bibliografia}

Les referències següents cobreixen tres funcions: introduccions al mateix nivell que aquest text, obres de consulta i fonts per a la transició cap a la lògica categòrica, i les fonts de la terminologia matemàtica catalana emprada.

\begin{thebibliography}{99}

\bibitem{awodey}
S.~Awodey,
\emph{Category Theory},
2a edició, Oxford Logic Guides~52, Oxford University Press, 2010.
\newblock (Text introductori de referència; guia principal per a l'ordre i el nivell d'aquest llibre, i per als capítols sobre categories cartesianes tancades i lògica.)

\bibitem{maclane}
S.~Mac~Lane,
\emph{Categories for the Working Mathematician},
2a edició, Graduate Texts in Mathematics~5, Springer, 1998.
\newblock (Obra de consulta clàssica sobre teoria de categories.)

\bibitem{leinster}
T.~Leinster,
\emph{Basic Category Theory},
Cambridge Studies in Advanced Mathematics~143, Cambridge University Press, 2014.
\newblock (Introducció alternativa, breu i moderna; disponible també en accés obert.)

\bibitem{riehl}
E.~Riehl,
\emph{Category Theory in Context},
Aurora: Dover Modern Math Originals, Dover, 2016.
\newblock (Consulta complementària, amb molts exemples.)

\bibitem{borceux}
F.~Borceux,
\emph{Handbook of Categorical Algebra},
3 volums, Encyclopedia of Mathematics and its Applications, Cambridge University Press, 1994.
\newblock (Referència exhaustiva per a límits, adjuncions i exactitud.)

\bibitem{maclanemoerdijk}
S.~Mac~Lane i I.~Moerdijk,
\emph{Sheaves in Geometry and Logic: A First Introduction to Topos Theory},
Universitext, Springer, 1994.
\newblock (Pont cap a la teoria de topos i la lògica geomètrica.)

\bibitem{lambekscott}
J.~Lambek i P.~J.~Scott,
\emph{Introduction to Higher-Order Categorical Logic},
Cambridge Studies in Advanced Mathematics~7, Cambridge University Press, 1986.
\newblock (Correspondència entre lògica d'ordre superior, càlcul lambda i categories cartesianes tancades.)

\bibitem{jacobs}
B.~Jacobs,
\emph{Categorical Logic and Type Theory},
Studies in Logic and the Foundations of Mathematics~141, Elsevier, 1999.
\newblock (Fibracions, hiperdoctrines i teoria de tipus des del punt de vista categòric.)

\bibitem{johnstone}
P.~T.~Johnstone,
\emph{Sketches of an Elephant: A Topos Theory Compendium},
Oxford Logic Guides 43--44, Oxford University Press, 2002.
\newblock (Obra avançada de consulta sobre teoria de topos.)

\bibitem{lawvere-ccfm}
F.~W.~Lawvere,
\emph{The Category of Categories as a Foundation for Mathematics},
en \emph{Proceedings of the Conference on Categorical Algebra (La Jolla, 1965)}, Springer, 1966, pàg.~1--20.
\newblock (Teoria elemental de categories en primer ordre; precedent de la presentació relacional de l'apèndix~\ref{cap:axiomatica-relacional}.)

\bibitem{scott-identity}
D.~S.~Scott,
\emph{Identity and Existence in Intuitionistic Logic},
en M.~Fourman, C.~Mulvey i D.~Scott (eds.), \emph{Applications of Sheaves}, Lecture Notes in Mathematics~753, Springer, 1979, pàg.~660--696.
\newblock (Operacions parcials, existència i identitat; rerefons de les axiomàtiques \enquote{només amb fletxes}.)

\bibitem{freydscedrov}
P.~J.~Freyd i A.~Scedrov,
\emph{Categories, Allegories},
North-Holland Mathematical Library~39, North-Holland, 1990.
\newblock (Presentació sistemàtica sense objectes primitius.)

\bibitem{BS2020}
C.~Benzmüller i D.~S.~Scott,
\emph{Automating Free Logic in HOL, with an Experimental Application in Category Theory},
Journal of Automated Reasoning~64, 2020, pàg.~53--72.
\newblock (Formalització en Isabelle/HOL de diverses axiomàtiques de categories amb lògica lliure.)

\bibitem{BSAFP}
C.~Benzmüller i D.~S.~Scott,
\emph{Axiom Systems for Category Theory in Free Logic},
Archive of Formal Proofs, 2018. \url{https://isa-afp.org/entries/AxiomaticCategoryTheory.html}.
\newblock (Comparació i verificació d'equivalències entre sistemes de Mac~Lane, Scott i Freyd--Scedrov.)

\bibitem{diec}
Institut d'Estudis Catalans,
\emph{Diccionari de la llengua catalana},
2a edició, Barcelona. \url{https://dlc.iec.cat}.
\newblock (Font normativa de la terminologia general catalana.)

\bibitem{termcat}
Universitat Politècnica de Catalunya; Enciclopèdia Catalana,
\emph{Diccionari de matemàtiques i estadística},
1a ed., Barcelona, 2002; 2a ed. en línia, Institut d'Estudis Catalans, 2021. \url{https://cit.iec.cat/DME/}.
\newblock (Font de la terminologia matemàtica catalana: functor, morfisme, preordre, etc.)

\end{thebibliography}
